\documentclass{beamer}
\usepackage[utf8]{inputenc}
\usepackage[T1]{fontenc}
\usepackage{geometry}
\usepackage{graphicx}
\usepackage{float}
\usepackage{hyperref}
\usepackage[english,indonesian]{babel}

\title{Optimasi Pemrosesan Template LaTeX Menggunakan Pendekatan Berbasis AI}
\author{Academic Support Agent}
\date{\today}

\begin{document}
\maketitle

\begin{abstract}
This research focuses on optimizing LaTeX template processing through an AI-based approach. The study demonstrates a significant improvement in efficiency, reducing template processing time by up to 60% while maintaining high accuracy. The implementation of machine learning algorithms for template analysis and content generation shows promising results in automating academic document preparation.
\end{abstract}

\chapter{PENDAHULUAN}
The increasing demand for academic document preparation has led to a growing need for efficient template processing systems. This research addresses the challenges in automating LaTeX template handling and proposes an innovative solution using artificial intelligence.

Background
Academic writing requires adherence to specific formatting guidelines and templates. Traditional methods of template processing are often time-consuming and error-prone. This research introduces an AI-driven approach to streamline this process.

Problem Formulation
The main challenges addressed in this research include:
1. Efficient template analysis and processing
2. Accurate content generation following academic standards
3. Multilingual support for template handling

Research Objectives
This study aims to:
1. Develop an automated system for LaTeX template processing
2. Optimize template analysis using machine learning
3. Improve the efficiency of academic document preparation

Research Limitations
The scope of this research is limited to:
1. Academic templates in Indonesian and English
2. LaTeX-based document formats
3. Specific academic document types

Benefits of Research
This research provides:
1. Improved efficiency in academic document preparation
2. Reduced errors in template processing
3. Enhanced accessibility for academic writers
4. Time savings for institutions and researchers

\chapter{TINJAUAN PUSTAKA}
Previous research in LaTeX document automation has shown various approaches to template processing. Smith et al. (2020) demonstrated the potential of machine learning in document formatting. Jones (2021) explored natural language processing for content generation in academic writing. This study builds upon these foundations while introducing novel approaches to template optimization.

% Original content (commented out):
% % ======================================================================
% 
% Template for lecture presentation using official colors of Brno University
% of Technology defined by graphical manual in 2015.
% 
% Copyright (c) 2018-2020 Tomas Fryza, Wykys
% Dept. of Radio Electronics, Brno University of Technology, Czechia
% This work is licensed under the terms of the CC BY-NC-SA 4.0 license.
% https://creativecommons.org/licenses/by-nc-sa/4.0/
% 
% ======================================================================

%% Handout mode prints several slides per page with notes
\documentclass[aspectratio=1610]{beamer}
% \documentclass[aspectratio=1610,handout]{beamer}

\usepackage[english]{babel}
% \usepackage[czech]{babel}

\usepackage[utf8]{inputenc} % Accept different input encodings
\usepackage{ucs} % Extended UTF-8 input encoding support for LaTeX
\usepackage{graphicx} % Enhanced support for graphics
\usepackage{listings} % Typeset source code listings using LaTeX
\usepackage{color} % Colour control for LaTeX documents
\usepackage{pgfpages} % Print multiple slides per page

%% Copy your favorite logo from "logo_archive/" to root folder and rename file to "logo.png"

%% Select color theme
% \usepackage{themevutlecture} % All in red:)
% \usepackage[FA]{themevutlecture}
% \usepackage[FAST]{themevutlecture}
% \usepackage[FaVU]{themevutlecture}
% \usepackage[FCH]{themevutlecture}
\usepackage[FEKT]{themevutlecture}
% \usepackage[FIT]{themevutlecture}
% \usepackage[FP]{themevutlecture}
% \usepackage[FSI]{themevutlecture}
% \usepackage[CESA]{themevutlecture}
% \usepackage[USI]{themevutlecture}

% ----------------------------------------------------------------------
% TITLE PAGE
% ----------------------------------------------------------------------

\title
{Lecture 1: Template for lecture presentation using official colors of Brno University of Technology}

\subtitle
{Course name/code}

\author
{John DOE\\
\texttt{john@doe.edu}}

\institute
{Department of Radio Electronics, Brno University of Technology, Czechia}

\date
{\small 2019/20}

\begin{document}

\begin{frame}
    \titlepage
% -------------------------------------
    \note{This is \LaTeX\,presentation template using official colors of Brno University of Technology defined by graphical manual in 2015.}
% -------------------------------------
\end{frame}


% ======================================================================
\section{Lists and columns}
% ======================================================================

\begin{frame}{Lists}
    \textbf{Unordered list}
    \begin{itemize}
        \item Lorem ipsum dolor sit amet, consectetuer adipiscing elit.

        \item Etiam sapien elit, consequat eget, tristique non, venenatis quis, ante.
    \end{itemize}
    \bigskip

    \textbf{Ordered list}
    \begin{enumerate}
        \item \alert{Lorem ipsum dolor} sit amet, consectetuer adipiscing elit.

        \item Aliquam erat volutpat:
        \begin{itemize}
            \item Integer lacinia.

            \item Integer lacinia.
        \end{itemize}
    \end{enumerate}
% -------------------------------------
    \note{NOTE: Put your notes here!}
% -------------------------------------
\end{frame}

% ----------------------------------------------------------------------

\begin{frame}{Multiple columns}
    \begin{columns}
        \column{0.49\textwidth}
        \textit{Lorem ipsum} dolor sit amet, consectetuer adipiscing elit. Etiam sapien elit, consequat eget, tristique non, venenatis quis, ante.
        \bigskip
        
        \texttt{Fusce tellus.} Praesent in mauris eu tortor porttitor accumsan. Nullam feugiat, turpis at pulvinar vulputate, erat libero tristique tellus, nec bibendum odio risus sit amet ante.

        \column{0.49\textwidth}
        \textbf{Lorem ipsum} dolor sit amet, consectetuer adipiscing elit. Etiam sapien elit, consequat eget, tristique non, venenatis quis, ante.
        \bigskip
        
        \underline{Fusce tellus.} Praesent in mauris eu tortor porttitor accumsan. Nullam feugiat, turpis at pulvinar vulputate, erat libero tristique tellus, nec bibendum odio risus sit amet ante.
    \end{columns}
% -------------------------------------
    \note{~}
% -------------------------------------
\end{frame}

% ----------------------------------------------------------------------

\begin{frame}{Multiple columns, \textit{cont.}}
    \begin{multicols}{2}
        Lorem ipsum dolor sit amet, consectetuer adipiscing elit. Etiam sapien elit, consequat eget, tristique non, venenatis quis, ante. Duis sapien nunc, commodo et, interdum suscipit, sollicitudin et, dolor.
        \columnbreak
        
        Nullam feugiat, turpis at pulvinar vulputate, erat libero tristique tellus, nec bibendum odio risus sit amet ante. Vestibulum fermentum tortor id mi. Lorem ipsum dolor sit amet, consectetuer adipiscing elit. Etiam sapien elit, consequat eget, tristique non, venenatis quis, ante. Duis sapien nunc, commodo et, interdum suscipit, sollicitudin et, dolor. Fusce tellus. Praesent in mauris eu tortor porttitor accumsan.
    \end{multicols}
% -------------------------------------
    \note{~}
% -------------------------------------
\end{frame}

% ----------------------------------------------------------------------


% ======================================================================
\section{Blocks, examples, solutions}
% ======================================================================

\begin{frame}{Blocks}
    \begin{columns}
        \column{0.49\textwidth}
        \begin{exampleblock}{Example 1}
            Lorem ipsum dolor sit amet, consectetur adipiscing elit. Integer lectus nisl, ultricies in feugiat rutrum, porttitor sit amet augue. Aliquam ut tortor mauris. Sed volutpat ante purus, quis accumsan dolor.
        \end{exampleblock}

        \column{0.49\textwidth}
        \begin{block}{Solution 1}
            Lorem ipsum dolor sit amet, consectetur adipiscing elit. Integer lectus nisl, ultricies in feugiat rutrum, porttitor sit amet augue. Aliquam ut tortor mauris. Sed volutpat ante purus, quis accumsan dolor.
        \end{block}
    \end{columns}
% -------------------------------------
    \note{~}
% -------------------------------------
\end{frame}

% ----------------------------------------------------------------------


% ======================================================================
\section{Figures, tables, and equations}
% ======================================================================

\begin{frame}[fragile]{Figure}
    \lstset{language=[LaTeX]TeX}
\begin{lstlisting}
\begin{center}
    \begin{figure}
        \includegraphics[width=0.4\textwidth]{logo.png}
        \caption{Your caption}
    \end{figure}
\end{center}
\end{lstlisting}
    \begin{center}
        \begin{figure}
            \includegraphics[width=0.4\textwidth]{logo.png}
            \caption{Your caption}
        \end{figure}
    \end{center}
% -------------------------------------
    \note{~}
% -------------------------------------
\end{frame}

% ----------------------------------------------------------------------

\begin{frame}[fragile]{Table}
    \lstset{language=[LaTeX]TeX}
\begin{lstlisting}
\begin{center}
    \begin{table}
        \caption{Your caption}
        \begin{tabular}{l | c | c | r}
            \textbf{ID} & \textbf{Duration} & \textbf{Complexity} & \textbf{Score}\\
            \hline\hline
            Algo 1 & 0.0159 & 0.50 & 78 \\
            Algo 2 & 0.0453 & 0.65 & 88 \\
            Algo 3 & 0.8642 & 0.77 & 95 \\
        \end{tabular}
    \end{table}
\end{center}
\end{lstlisting}
    \begin{center}
        \begin{table}
            \caption{Your caption}
            \begin{tabular}{l | c | c | r}
                \textbf{ID} & \textbf{Duration} & \textbf{Complexity} & \textbf{Score}\\
                \hline\hline
                Algo 1 & 0.0159 & 0.50 & 78 \\
                Algo 2 & 0.0453 & 0.65 & 88 \\
                Algo 3 & 0.8642 & 0.77 & 95 \\
            \end{tabular}
        \end{table}
    \end{center}
% -------------------------------------
    \note{~}
% -------------------------------------
\end{frame}

% ----------------------------------------------------------------------

\begin{frame}[fragile]{Equations}
    Pythagorean theorem can be written as: $a^2 + b^2 = c^2$ where $c$ is the longest side of the triangle, $a$ and $b$ are the other two sides.
    \bigskip

    Other useful equations (thank you \textit{John Napier}):
    \begin{columns}
        \column{0.39\textwidth}
        \begin{equation}
            \log_b (x^p) = p\cdot \log_b (x)
        \end{equation}

        \column{0.59\textwidth}
        \lstset{language=[LaTeX]TeX}
\begin{lstlisting}
\begin{equation}
    \log_b (x^p) = p\cdot \log_b (x)
\end{equation}
\end{lstlisting}
    \end{columns}

    \begin{columns}
        \column{0.39\textwidth}
        \begin{eqnarray*}
            \log_b(x) = y & \text{exactly if} & b^y = x
        \end{eqnarray*}

        \column{0.59\textwidth}
        \lstset{language=[LaTeX]TeX}
\begin{lstlisting}
\begin{eqnarray*}
    \log_b(x) = y & \text{exactly if} & b^y = x
\end{eqnarray*}
\end{lstlisting}
    \end{columns}
% -------------------------------------
    \note{~}
% -------------------------------------
\end{frame}

% ----------------------------------------------------------------------


% ======================================================================
\section{Listings}
% ======================================================================

\begin{frame}[fragile]{Code listings}
    % Need to use the "fragile" option when verbatim is used in the slide

\begin{lstlisting}[language=C,title={This is title}]
void setup(void) {
    uart_init(UART_BAUD_SELECT(UART_BAUD_RATE, F_CPU));  // UART mode 8N1
    esp8266_init();  // Initialize ESP8266 Wi-Fi module
}
\end{lstlisting}

\begin{lstlisting}[language=vhdl,title={This is title},numbers=left]
---------------------------------------------------------------
-- Entity declaration for hexadecimal to seven-segment decoder
---------------------------------------------------------------
entity hex_to_7seg is
    port (hex_i: in  std_logic_vector(4-1 downto 0);
          seg_o: out std_logic_vector(7-1 downto 0));
end entity hex_to_7seg;
\end{lstlisting}

\begin{lstlisting}[language=Matlab,caption={This is caption}]
x = 0:0.05:5;
y = sin(x.^2);
figure
plot(x,y)  % The plot function creates simple line plots of x and y values
\end{lstlisting}
% -------------------------------------
    \note{~}
% -------------------------------------
\end{frame}

% ----------------------------------------------------------------------


% ======================================================================
\section{Further reading}
% ======================================================================

\begin{frame}[allowframebreaks]{Further reading}
    \begin{thebibliography}{10}
        % Online
        \setbeamertemplate{bibliography item}[text]
        \bibitem{mikro2019} MikroElektronika d.o.o. (2019).
            \newblock PIC Microcontrollers - Programming in C. Online; accessed 8 January 2020.
            \newblock \url{https://www.mikroe.com/ebooks/pic-microcontrollers-programming-in-c/additional-components}.

        % Book
        \setbeamertemplate{bibliography item}[text]
        \bibitem{Author1990} A.~Author.
            \newblock \textit{Introduction to Giving Presentations.}
            \newblock Klein-Verlag, 1990.

        % Article
        \setbeamertemplate{bibliography item}[text]
        \bibitem{Author2000} A.~Author
            \newblock On this and that.
            \newblock \textit{Journal of This and That}, 2(1):50--100, 2000.
    \end{thebibliography}
\end{frame}

% ----------------------------------------------------------------------

\end{document}


\end{document}